\section{Negative branch and the scalar side of the problem}
\label{sec:negative}
% Status: derived/conjectural. Branch algebra is exact; scalar interpretation requires a dynamical map.

The negative branch in Eq.~\eqref{eq:roots} is fixed by the same secular equation as the positive branch.  It is therefore part of the clue.  Hans de Vries identified the positive-branch proximity to the mass-shell weak angle.  Under the temporary normalization \(M_Z^2=\mu^2x_+(2)\), the negative magnitudes land near electroweak scalar and vacuum-ray scales.  That arithmetic observation supplies the main source-selection question: can one source theory attach \(x_-(J)\) to a gauge-invariant scalar functional generated by the same reduced operator that gives \(x_+(J)\)?

The root identities imply
\begin{equation}
  x_-(J)=-J-x_+(J)=-\frac{J}{x_+(J)}.
  \label{eq:negative-from-positive}
\end{equation}
Any physical use of the positive branch automatically carries a partner branch.

The Higgs potential can be written as
\begin{equation}
  V(H)=m^2 H^\dagger H+\lambda (H^\dagger H)^2,
  \qquad m^2<0
  \label{eq:higgs-potential}
\end{equation}
in the broken phase.  The sign of the quadratic term makes the scalar interpretation a concrete target.  If the normalization is fixed by \(M_Z^2=\mu^2 x_+(2)\), the two descriptive branch magnitudes are
\begin{equation}
  M_{-,1} = M_Z\sqrt{\frac{|x_-(2)|}{x_+(2)}},
  \qquad
  M_{-,1/2} = M_Z\sqrt{\frac{|x_-(3/4)|}{x_+(2)}} .
  \label{eq:negative-scales}
\end{equation}
Under this normalization, the two branch magnitudes are descriptive arithmetic.  They become physical scalar identifications after the same determinant, normalization, selected \(J_\star\), and scheme are derived from gauge-Higgs or source-theory data.

A valid identification must be gauge-invariant and scheme-controlled.  Possible targets include \(m^2\), \(\lambda v^2\), the Higgs pole mass, a Wilson-line curvature, or an eigenvalue in a compactification problem.  The branch values associated with \(J=3/4\) and \(J=2\) generate definite dimensionless numbers.  The map from those numbers to scalar observables remains open.

\subsection{Logical role of the negative root}
\label{sec:negative-logical-role}

The negative root changes the character of the clue.  The full quadratic gives a paired object: a pole-ratio candidate and a scalar-side candidate arise from one secular equation.  This pairing supplies the main conceptual reason to retain the negative branch.  A dynamical theory that produces Eq.~\eqref{eq:intro-secular} has to say what both branches mean in the same basis.

The two-root structure can be stated as a spectral chain,
\begin{equation}
  \mathcal M_J
  \longrightarrow
  \left\{
  \begin{array}{ll}
  x_+(J) & \hbox{projective vector datum},\\[2pt]
  x_-(J) & \hbox{candidate scalar/order-parameter target}.
  \end{array}
  \right.
  \label{eq:negative-chain}
\end{equation}
The first line is the pole-ratio reading developed in Sec.~\ref{sec:pole}.  The second line is a target for the Higgs, boundary, brane, or compactification sector.  The chain is meaningful because both entries come from the same diagonalization.  Any proposed derivation that keeps the positive branch acquires the negative-branch obligation automatically.

This obligation also sharpens the electroweak assignment.  The ordered pair \((J_H,J_{\rm adj})=(3/4,2)\) already suggests two electroweak structures: the doublet order parameter and the adjoint current.  The negative branch asks whether the doublet/order-parameter side can be read as a scalar instability, a vacuum scale, or a deformation mode in the same operator that gives the vector pole quotient.  The scalar reading therefore ties together three open calculations: the ordered \(J\)-assignment, the determinant derivation, and the branch-to-Higgs map.

\subsection{Gauge-invariant scalar language}
\label{sec:gauge-invariant-scalar-language}

The scalar sector has several layers.  The Higgs doublet \(H\) itself is gauge-covariant.  The radial composite
\begin{equation}
  \varphi^2 = 2H^\dagger H
  \label{eq:radial-composite}
\end{equation}
is the local invariant order-parameter language for this manuscript's purpose after renormalization and scheme choice.  In the broken phase, a gauge choice may represent the vacuum by a real radial expectation value, yet the physical statement should be expressible through invariant quantities, pole positions, or matched parameters.  The negative branch therefore requires a defined scalar functional before it can carry a physical claim.

An abstract target is
\begin{equation}
  \mathcal F_{\rm sc}
  =
  \mathcal F_{\rm sc}
  \big(
  H^\dagger H,\,
  (D_\mu H)^\dagger D^\mu H,\,
  V(H),\,
  \hbox{boundary or compactification data}
  \big),
  \label{eq:scalar-functional-general}
\end{equation}
with a branch assignment
\begin{equation}
  \mathcal F_{\rm sc}(J_\star;\Lambda,s_{\rm ren},u)
  =
  C_{\rm sc}(\Lambda,s_{\rm ren},u)\,|x_-(J_\star)|.
  \label{eq:scalar-functional-branch}
\end{equation}
Here \(J_\star\), \(C_{\rm sc}\), \(\Lambda\), the renormalization scheme \(s_{\rm ren}\), and the source datum \(u\) are part of the theorem.  The expression may describe a pole mass, a curvature of the scalar potential, a vacuum-ray coordinate, a brane separation, a Wilson-line curvature, or a compactification deformation.  The manuscript treats all of these as candidate targets until one route derives a specific object.

The same-source package needed for a physical reading is
\begin{equation}
  u\longmapsto
  \big(
  \mathcal H_J,\,
  \langle\cdot,\cdot\rangle_u,\,
  P_J,\,
  \Lambda_J,\,
  \mathcal R_{\rm pole},\,
  \mathcal F_{\rm sc}
  \big),
  \label{eq:negative-same-source-package}
\end{equation}
with \(\mathcal H_J\) the reduced two-channel space, \(P_J\) the projection,
\(\Lambda_J\) the source scale, and \(\mathcal R_{\rm pole}\) the pole map used
for the vector branch.  A dimensionless scalar target then has the form
\begin{equation}
  \widehat{\mathcal F}_{\rm sc}(u;J_\star)
  =
  \Lambda_J^{-d_{\rm sc}}\,
  \mathcal F_{\rm sc}(J_\star;\Lambda_J,s_{\rm ren},u)
  =
  C_{\rm sc}(s_{\rm ren},u)\,|x_-(J_\star)|,
  \label{eq:negative-hatted-scalar-functional}
\end{equation}
where \(d_{\rm sc}\) is the mass dimension of the chosen observable.  The source
derivation must state whether \(\widehat{\mathcal F}_{\rm sc}\) is a vacuum-ray
coordinate, a potential curvature, a pole observable, a boundary modulus, or a
compactification eigenvalue.

The target can be summarized by the commutative source diagram
\begin{equation}
\begin{array}{ccccc}
\mathcal L_{\rm src}(u)
&\longrightarrow&
(\mathcal H_J,\langle\cdot,\cdot\rangle_u,P_J)
&\longrightarrow&
Q(J)
\\
&&&&\downarrow
\\
&&&&
\left\{
\begin{array}{ll}
x_+(J) & \hbox{vector pole quotient},\\[2pt]
x_-(J) & \hbox{candidate scalar functional},\\[2pt]
x_-(J) & \hbox{auxiliary or gauge-fixed branch}.
\end{array}
\right.
\end{array}
\label{eq:negative-source-diagram}
\end{equation}
The third line is the explicit failure outcome for a scalar interpretation.

A sharper downstream target appears if the scalar functional is the
vacuum-ray normalization.  In that case the Standard Model identity
\(M_W^2=g^2v^2/4\) can turn the pair \(x_+(3/4)\) and \(x_-(2)\) into a
candidate construction coupling.  Appendix~\ref{app:target-alpha-endpoint}
records this as an electromagnetic endpoint theorem target: the scalar branch,
the charged-vector normalization, the \(D=9\) \(U(1)_{\rm em}\) kinetic term,
and the matching pair \((Q_\alpha,\Delta_{\alpha,{\rm match}})\) must be
derived before \(\alpha\) acquires construction status.

The scalar pole target is the most restrictive because it demands the same pole discipline as the vector comparison.  The potential target is more sensitive to matching conventions, because \(m^2\), \(\lambda\), and \(v\) depend on the renormalization scheme and on how the effective potential is defined.  The compactification target is broader: \(x_-(J)\) may first appear as an eigenvalue or deformation parameter in a higher-dimensional problem, with the four-dimensional Higgs interpretation emerging after reduction.

\subsection{Vacuum ray and branch pairing}
\label{sec:vacuum-ray-branch-pairing}

The electroweak ray separates scale from direction.  Along the broken-to-unbroken ray, the radial coordinate scales the vector masses, while the ratio
\begin{equation}
  \frac{M_W^2}{M_Z^2}
  =
  \frac{g^2}{g^2+g'^2}
  \label{eq:negative-vector-projector}
\end{equation}
is projective at tree level.  The DeVries quotient addresses this projective datum.  The negative branch is a scalar-side target because it shares the same two-channel operator while carrying the opposite eigenvalue.

In a route with fields \(h_J\) and \(a_J\), the branch vectors have the schematic form
\begin{equation}
  u_{\pm}(J)=\alpha_\pm(J)h_J+\beta_\pm(J)a_J.
  \label{eq:negative-branch-mixture}
\end{equation}
The scalar-side reading becomes a statement about the \(h_J\) component of \(u_-(J)\).  The vector-side reading becomes a statement about the projective pole ratio extracted from \(u_+(J)\).  A complete derivation should therefore give a branch basis, an inner product, and a normalization.  These data determine whether the negative eigenvector genuinely selects an order-parameter amplitude or whether it belongs to an auxiliary sector of the chosen source theory.

The electroweak ray also supplies a consistency check.  Taking the radial vacuum coordinate toward the symmetric point should collapse the vector masses together with the Higgs order parameter.  The DeVries ratio can remain as a projective direction in coupling or spectral space.  A scalar interpretation of \(x_-\) should respect this ray structure by scaling with the radial data or by controlling the deformation that sets the radial scale.

\subsection{Candidate scalar targets}

The Standard Model scalar sector gives several inequivalent objects that can carry a scalar-side interpretation \cite{Logan2022,Dawson2018}.  They have different scheme and gauge properties:
\begin{center}
\small
\begin{tabular}{p{0.24\linewidth}p{0.34\linewidth}p{0.32\linewidth}}
\toprule
Target & Meaning & DeVries requirement\\
\midrule
\(m^2(s_{\rm ren},\Lambda)\) & quadratic term in the Higgs potential & define the scheme and invariant extraction\\
\(\lambda(s_{\rm ren},\Lambda)v^2\) & curvature scale along the radial mode & connect branch value to the electroweak ray\\
\(M_{h,\rm pole}^2\) & scalar pole position & derive the scalar pole from the same determinant as W/Z\\
\(\partial_\alpha^2V_{\rm eff}(\alpha)|_{\alpha_\star}\) & Wilson-line or Hosotani curvature & identify the holonomy and source normalization\\
radion/modulus eigenvalue & compactification scalar or boundary deformation & identify the compact operator and its normalization\\
brane separation mode & endpoint or brane order parameter & show coupling to the adjoint/current channel\\
\bottomrule
\end{tabular}
\end{center}

The pole interpretation is the sharpest target.  It asks the same object that governs the vector pole quotient to generate a scalar pole or scalar instability.  The potential interpretation is more flexible, because \(m^2\), \(\lambda\), and \(v\) depend on conventions and matching.  The compactification interpretation is broader still; it may identify the negative branch with a modulus or boundary deformation whose four-dimensional scalar meaning emerges after reduction.

\subsection{Route-specific scalar maps}
\label{sec:route-specific-scalar-maps}

The endpoint route maps the negative branch to brane scalar data.  In a D-brane system, transverse positions of branes appear as scalar fields on the worldvolume, while strings stretched between branes acquire masses controlled by separation data \cite{TongString,Polchinski1994}.  A DeVries endpoint derivation would identify \(u_-(J)\) with a brane-Higgsing direction, a stretched-string scalar, or a boundary order-parameter field.  The map has to use the same endpoint algebra that supplies the current trace and off-diagonal product.

The interval route maps the negative branch to boundary or modulus data.  A five-dimensional gauge field supplies four-dimensional vector modes and scalar components; the compactification radius, brane terms, and boundary conditions determine the light spectrum \cite{CsakiHubiszMeade2005,WittenKK1981}.  A DeVries interval derivation would identify \(u_-(J)\) with a boundary scalar amplitude, an \(A_5\)-like scalar, a radion-like deformation, or a brane-localized order parameter.  The map has to preserve the same boundary determinant that gives the W/Z branch quotient.  The Hosotani version of this target uses the holonomy variable
\begin{equation}
  \alpha_H \sim g_5\int_I dy\,A_5 ,
  \qquad
  \mathcal F_{\rm sc}^{\rm Hol}
  =
  \left.
  \frac{\partial^2 V_{\rm eff}(\alpha_H;\Lambda,s_{\rm ren})}
       {\partial\alpha_H^2}
  \right|_{\alpha_H=\alpha_\star},
  \label{eq:hosotani-scalar-functional}
\end{equation}
and asks whether the same interval kernel that yields the transverse vector pole
also yields
\(\widehat{\mathcal F}_{\rm sc}^{\rm Hol}
=C_{\rm sc}|x_-(J_\star)|\) \cite{Hosotani1983,HabaOda2011,Bucci2004}.

The \(G_2\) route maps the negative branch to singularity deformation data.  Local \(G_2\) constructions place gauge sectors and charged matter at singular loci, and Higgs-bundle descriptions encode matter localization and interaction data through functions on an associative three-cycle \cite{AcharyaWitten2001,AcharyaGukov2004,BraunEtAl2019}.  A DeVries \(G_2\) derivation would identify \(u_-(J)\) with a deformation parameter, a localized modulus, a critical value of \(f_Q\), or a charge-density functional.  The same local geometry then has to produce the adjoint/current and order-parameter entries of the kernel.

The route ledger is:
\begin{center}
\small
\begin{tabular}{p{0.19\linewidth}p{0.35\linewidth}p{0.36\linewidth}}
\toprule
Route & Source-supported scalar variable & Missing DeVries map\\
\midrule
Endpoint & transverse brane scalar, stretched-string separation, boundary order parameter & derive the same endpoint determinant, projection, and normalization for \(u_-(J)\)\\
Interval & boundary Higgs datum, \(A_5\) holonomy, radion, boundary-condition modulus & derive the same interval kernel, gauge fixing, scalar/eaten separation, and \(\Lambda_J\)\\
\(G_2\) & singularity deformation, Higgs-bundle one-form, localized modulus & derive the local two-channel kernel and the scalar functional on the same associative-cycle data\\
\bottomrule
\end{tabular}
\end{center}

These maps give the negative branch a practical role in route selection.  A candidate route that produces the positive-branch quotient and gives the negative eigenvector a scalar or deformation interpretation gains additional physical content.  A route that assigns the negative branch to an auxiliary variable defines a narrower physical claim.

\subsection{Eigenvectors and branch meaning}

The two roots also have eigenvectors.  For the matrix in Eq.~\eqref{eq:QJ}, an eigenvector with eigenvalue \(x\) can be represented by
\begin{equation}
  v_x \propto
  \begin{pmatrix}
  \sqrt J\\
  x
  \end{pmatrix}.
  \label{eq:branch-eigenvector}
\end{equation}
Thus the branch assignment is a statement about channel composition.  The positive and negative roots are mixtures of the same two basis amplitudes.  A scalar interpretation of \(x_-\) should therefore identify the basis in which one component is the vector/current side and the other is the order-parameter or boundary side.

A criterion is the following.  If the same two-channel determinant produces both the vector pole ratio and the scalar instability, then the determinant should appear in an effective action, boundary condition, or self-energy equation.  The corresponding theorem target is an explicit determinant with an identified second-channel field or mode.  Breitenlohner--Freedman stability supplies a route caveat for AdS scalar readings: a negative mass-squared can be admissible under specified boundary conditions, while the DeVries sign still requires the source functional in Eq.~\eqref{eq:negative-hatted-scalar-functional} \cite{BreitenlohnerFreedman1982}.
